\section{Experiments}
\label{sec:experiments}
% In our experiments, we first examine both the benefits of an interactive setting towards coding tasks along with the gaps and research challenges it presents that models have yet to overcome.
% We then demonstrate how the interactive coding task typically encapsulates a multitude of reasoning challenges, motivating a higher order question of how models might go about exercising and alternating across multiple problem solving strategies on its own volition.
% Finally, we present a small-scale dataset inspired by Capture the Flag, a popular genre of computer security competitions, to showcase \benchmark{}'s value as a testbed for new datasets and task formats in the realm of interactive code.
% Specific trajectories that are evidence of the conclusions developed across the following sections can be found in \S~\ref{appx:experiments:trajectory_analyses}.

\begin{table}
\centering
\begin{tabular}{llllll|lllll}
\toprule
      \benchmark{}-SQL & \multicolumn{5}{c}{Single Turn} & \multicolumn{5}{c}{Try Again ($n$ = 10)} \\
     Model / Hardness & Easy & Med & Hard & Extra & All & Easy & Med & Hard & Extra & All \\
\midrule
text-davinci-003 & 20.6 & 4.9 & 1.7 & 0.0 & 7.4 & 32.4 & 14.6 & 5.2 & 4.2 & 15.6 \\
gpt-3.5-turbo & 22.6 & 8.3 & \textbf{5.7} & \textbf{3.6} & 10.5 & 72.5 & 44.3 & 43.7 & 21.1 & 47.3 \\
gpt-4 & 19.8 & 7.2 & 4.6 & 3.0 & 9.1 & \textbf{87.5} & \textbf{76.7} & \textbf{66.7} & \textbf{52.4} & \textbf{73.7} \\
text-bison-001 & \textbf{23.8} & \textbf{10.9} & \textbf{5.7} & 0.6 & \textbf{11.5} & 27.0 & 12.3 & 5.7 & 0.6 & 12.9 \\
chat-bison-001 & 18.5 & 6.5 & 4.0 & 0.0 & 7.9 & 22.2 & 7.8 & 6.9 & 0.0 & 9.9 \\
% Claude-v1 & & & & & & & & & & \\
Vicuna-13B & 8.1 & 1.3 & 0.6 & 0.0 & 2.6 & 18.9 & 3.4 & 1.7 & 0.0 & 6.3 \\
StarChat-16B & 21.8 & 7.4 & 2.9 & 0.0 & 8.9 & 22.3 & 8.5 & 2.9 & 1.2 & 9.7 \\
\bottomrule
\end{tabular}
\vspace{5pt}
\caption{Success Rate for single vs.\,multi turn evaluation on \benchmark{}-SQL (refer \S\ref{appx:env:sqlenv}). Query difficulty is adopted from Spider~\cite{yu-etal-2018-spider}. Best metrics are in \textbf{bold}.}
\label{table:single-vs-multi-sql}
\vspace{-10pt}
\end{table}
\begin{table}
\centering
\begin{tabular}{llllll|lllll}
\toprule
     \benchmark{}-Bash & \multicolumn{5}{c}{Single Turn} & \multicolumn{5}{c}{Try Again ($n$ = 10)} \\
Model / File System & 1 & 2 & 3 & 4 & All & 1 & 2 & 3 & 4 & All \\
\midrule
text-davinci-003 & 10.0 & 32.1 & 28.8 & 33.3 & 24.6 & 30.0 & \textbf{52.8} & 32.2 & 44.4 & 38.7 \\
gpt-3.5-turbo & \textbf{30.0} & \textbf{39.6} & 33.3 & 37.0 & \textbf{34.5} & \textbf{45.0} & 49.1 & 45.0 & 48.1 & 46.5 \\
gpt-4 & 25.0 & 37.7 & \textbf{36.7} & \textbf{40.7} & 34.0 & 41.7 & 47.2 & \textbf{51.7} & \textbf{59.2} & \textbf{48.5} \\
text-bison-001 & 15.0 & 22.6 & 11.7 & 22.2 & 17.0 & 23.3 & 28.3 & 16.7 & 22.2 & 22.5 \\
chat-bison-001 & 12.1 & 22.5 & 16.7 & 22.2 & 17.7 & 13.8 & 24.5 & 18.3 & 22.2 & 19.2 \\
% Claude-v1 & & & & & & & & \\
Vicuna-13B & 10.0 & 24.5 & 18.3 & 7.4 & 16.0 & 15.0 & 35.8 & 25.0  & 22.2 & 24.5 \\
StarChat-16B & 15.5 & 22.6 & 13.3 & 22.2 & 17.7 & 17.2 & 30.2 & 21.7 & 29.6 & 23.7 \\
\bottomrule
\end{tabular}
\vspace{5pt}
\caption{Success Rate across file systems for single vs. multi-turn evaluation on  \benchmark{}-Bash (refer \S\ref{appx:env:bashenv}). To evaluate models' ability to interact with different task settings, we evaluate disjoint sets of Bash instructions across four different file systems. Best metrics are in \textbf{bold}.}
\label{table:single-vs-multi-bash}
\vspace{-1.5em}
\end{table}

\subsection{Base models comparison}
\label{sec:experiments:base_compare}
\textbf{Task performances.} We first compare the success rate of models in the Single Turn and Try Again settings for both the \benchmark{}-Bash and SQL datasets.
From Table~\ref{table:single-vs-multi-sql} and Table~\ref{table:single-vs-multi-bash}, we observe that performance across different levels of task difficulty (SQL) and different file systems (Bash) is superior in the interactive setting for all models, with a notable multi-fold increase for GPT-4 ($9.1\% \rightarrow 73.7\%$) on the \benchmark{}-SQL task.

\textbf{Analysis of interactions.} Manual inspection of trajectory logs indicates that models actively exercise later turns for discovering relevant context, correcting errors via execution feedback as observations, and solving problems via iteratively constructing and editing actions as affirmed by Figure~\ref{fig:turn_stats}.
In addition, models also demonstrate a level of planning and modular problem solving; for instructions with \texttt{gold} commands that chain multiple commands together (i.e. with \text{\textbar}, $>$, or \texttt{;} in \texttt{bash}) or consist of multiple sub-problems (i.e. subqueries in \texttt{SQL}), models will use observations from solving smaller sub-problems in earlier turns to compose the higher-order action. Trajectories that exhibit these phenomena are in \S~\ref{appx:experiments:trajectory_analyses}

\textbf{Failure cases.} With that said, both Figure~\ref{fig:turn_stats} exhibits a plateauing in Success Rate and and Error \%.
This suggests that as the amount of context and feedback builds up, models are less capable of discerning relevant past history toward future actions.
In late-turn scenarios, task episode trajectories often reveal repetition of earlier actions, a failure to effectively use recent observations towards deciding an appropriate next action, or an inability to recognize that a current problem-solving chain of thought is inconclusive or futile.
This is particularly evident for \texttt{hard} and \texttt{extra} level \benchmark{}-SQL task instructions that require context spanning across several tables and actions that incorporate multiple clauses.
We note that even when the full schema of all tables and their descriptions are offered in addition to the original instructions, models still benefit greatly from using interaction to experiment with different \texttt{JOIN} and filtering operators across multiple turns, as demonstrated in \S~\ref{appx:experiments:additional_expr}.
A larger context window size, retrieval of useful memory, and more adaptive reasoning paradigms are just a handful of potential solutions to overcoming such challenges.

\subsection{Prompting strategy comparison}
\label{sec:experiments:strategy_compare}
Initiating language agents with prompting strategies that encourage different forms of reasoning toward problem-solving improves performance on the interactive coding task to varying degrees. Table~\ref{table:strategy-comparison} presents side-by-side comparisons of the success rate, number of turns, and error rate per strategy. 
Compared to Try Again, which lacks specific guidance on leveraging multiple turns, more explicit reasoning frameworks such as ReAct and Plan \& Solve policies generally achieve higher success rates (SQL: $47.3\% \rightarrow 58.7\%$) with fewer turns and a higher rate of admissible commands.

\begin{figure}
    \begin{subfigure}[t]{0.5\textwidth}
        \centering
        \includegraphics[width=\textwidth]{figures/assets/bash_turn_stats.pdf}
        \caption{Success rate vs. turns for \benchmark{}-Bash}
    \end{subfigure}
     \begin{subfigure}[t]{0.5\textwidth}
     \centering
        \includegraphics[width=\textwidth]{figures/assets/sql_turn_stats.pdf}
        \caption{Success rate vs. turns for \benchmark{}-SQL}
    \end{subfigure}
    \caption{Growth in Success Rate with increase in number of interaction turns across models configured with Try Again prompting strategy for \benchmark{}-Bash and SQL tasks.}
    \label{fig:turn_stats}
\end{figure}
\begin{table}
\centering
\begin{tabular}{llll|lll|lll}
\toprule
     & \multicolumn{3}{c}{Try Again ($n$ = 10)} & \multicolumn{3}{c}{ReAct ($n$ = 10)} & \multicolumn{3}{c}{Plan \& Solve} \\
     & SR & Turns & Error \%
     & SR & Turns & Error \%
     & SR & Turns & Error \% \\
\midrule
SQL & 47.3 & 7.25 & 46.4
    & \textbf{58.7} & 5.30 & \textbf{6.94}
    & 49.1 & \textbf{4.29} & 16.2 \\
Bash & \textbf{46.5} & 6.15 & 24.9
    & 20.5 & \textbf{4.40} & \textbf{20.4}
    & 28.0 & 6.65 & 53.3 \\
\bottomrule
\end{tabular}
\vspace{5pt}
\caption{Comparison of different prompting strategies across the entire \benchmark{}-SQL and \benchmark{}-Bash datasets using \texttt{gpt-3.5-turbo} as the base model. \textit{Turns} refers to the average number of turns taken for a single task episode. For Try Again and ReAct, the max number of turns $n = 10$. The highest Success Rate, fewest Turns, and lowest Error \% are highlighted per dataset since they reflect more accuracy and efficient task solving. Best metrics are in \textbf{bold}.}
\label{table:strategy-comparison}
\vspace{-1.5em}
\end{table}

\textbf{Different tasks present different learning challenges.}
An important skill to solving the \benchmark{}-SQL task is the ability to discover context and construct actions conditionally based on information revealed in prior observations.
Given that \benchmark{}-SQL task instructions are phrased most commonly as questions, adapting to the task setting and new information discovered along the way puts more emphasis on error correction and context discovery.
On the other hand, the more declarative and multi-step nature of the \benchmark{}-Bash task instructions is more aptly solved by planning and modular task completion.
These distinctions manifest in the Plan \& Solve strategy's performance gap between the \benchmark{}-SQL and \benchmark{}-Bash tasks; while Plan \& Solve encourages a model to decompose problems into more manageable steps, the strategy is less favorable towards adjusting on the fly in response to execution feedback.
Example trajectories supporting these claims are in \S~\ref{appx:experiments:trajectory_analyses}.

\textbf{More adaptive reasoning is favorable.}
Compared to "imperative" reasoning paradigms such as Plan \& Solve which prescribe a relatively rigid procedure, more flexible frameworks like ReAct, which do not enforce any particular logical formula or roadmap, are more conducive to eliciting a broader set of reasoning capabilities.
However, while ReAct's performance is generally superior to Plan \& Solve, tasks solved by \textit{both} strategies with \texttt{gpt-3.5-turbo} make up $57\%$ ($407/708$) and $27.6\%$ ($21/76$) of the union of all successfully solved \benchmark{}-SQL and \benchmark{}-Bash tasks respectively.
This discrepancy highlights a trade-off between the guidance and structural constraints that are inherent to prompting strategies; schemes that draw out specific reasoning patterns often overlook other equally useful capabilities.
\benchmark{}'s interactive coding task can serve as a strong litmus test toward more adaptable, variegated model reasoning.

\subsection{New tasks \& datasets opportunities}
\benchmark{}'s task formulation, modular design, flexible task construction, and use of virtual containers enable task designers to manifest new, complex, code-driven tasks, where completion is much more attainable through interaction.
We draw inspiration from Capture the Flag (CTF)~\cite{cowan2003ctf}, a competitive cybersecurity game that requires expertise in coding, cryptography (i.e. binary exploitation, forensics), reverse engineering, and recognizing security vulnerabilities to accomplish the primary objective of discovering encrypted "flags" concealed within code snippets or file systems.
Compared to \benchmark{}-Bash \& -SQL, CTF is much more complicated, requiring an agent to exercise knowledge of multiple coding languages, modularize a higher-order objective into sub-problems, construct multi-step plans towards solving each problem, and adjust strategy when a plan fails to yield any useful insights.

We establish \benchmark{}-CTF, a new dataset consisting of 100 CTF objectives from picoCTF~\cite{picoctf}.
Following the interactive coding task formulation, each task instance in \benchmark{}-CTF is given as a \texttt{<instruction, assets, hidden flag>} tuple.
We first construct a Bourne Shell within an Ubuntu OS as the task environment.
Here, \benchmark{}'s use of virtual containers is crucial, as necessary actions can be irreversibly damaging on real systems (i.e. \texttt{rm -rf}, \texttt{sudo} access).
Per task instance, the associated assets (e.g., images, executables, code), necessary for task completion, are copied into the OS file system.
Given this setting, a task worker must understand the given material and investigate the assets to develop potential solutions.
Executing a successful approach must be done across multiple steps with various conditionals, where the execution feedback of a prior step could have a significant effect on the next step.
% For more complex problems, some solution result in an inconclusive dead end.
% In these situations, models must not only recognize when they should pivot their approach, but also decide what approach makes the most sense to attempt given prior failures.
Figure~\ref{fig:traj_ctf} spotlights the diverse skills needed for CTF. 

\begin{figure}
    \centering
    \makebox[\textwidth][c]{\includegraphics[width=\textwidth]{figures/assets/Traj-ctf-forensics-wide.pdf}}
    \caption{GPT-4's interaction trajectory for a binary exploitation CTF task. This requires proficiency in \textcolor{blue}{Bash} and \textcolor{OliveGreen}{Python}, among additional knowledge and reasoning. \textcolor{orange}{Orange} text and arrows highlight the feedback that the model attends to in generating the next action. In last step, agent submits flag.}
    \label{fig:traj_ctf}
\end{figure}

%%%% MARK: NON INCLUDED SECTIONS %%%%

% \subsection{Intercode Tasks Present Opportunities to design better techniques for soliciting different/multiple capabilities from LLMs}
% \jy{This message might not be worth explaining or the amount of explanation it would require (1. explain plan\&solve + refine and why it was made, 2. Discuss how it performs better than other strategies 3. Still not perfect)}
% \input{figures/tables/expr-plan-solve-refine}

% \subsection{Tools / Sim-to-Real Use Cases}
% Nothing actually done in this area, but is this a worthwhile thing to have + can we do something in the next two weeks? Can be a toy thing.

% IDEA: Take existing code and refactor it to be (more readable/efficient, handle some errant edge case). We don't have unit tests - use an environment to interactively verify that the refactored version works the same as before + achieves the improvement. (Requires Web + Bash/SQL?)